\DocumentMetadata{
  pdfversion=2.0,
  pdfstandard=ua-2,
  lang=en-US,
  tagging=on,
  tagging-setup={math/setup=mathml-SE,tabsorder=structure}
}
\documentclass[11pt,letterpaper]{article}
\usepackage[margin=0.68in,includefoot]{geometry}
\usepackage{newcomputermodern}
\usepackage{amsmath,amscd,mathtools}
\providecommand{\square}{\mdlgwhtsquare}
\usepackage[hidelinks]{hyperref}
\hypersetup{
  pdftitle={Probability PhD Exam, May 2023},
  pdfauthor={University of Florida Department of Mathematics},
  pdfsubject={Graduate mathematics examination},
  pdfdisplaydoctitle=true,
  bookmarksopen=true
}
\setcounter{secnumdepth}{0}
\setlength{\parindent}{0pt}
\setlength{\parskip}{0.28em}
\setlength{\emergencystretch}{3em}
\setlength{\leftmargini}{2.15em}
\setlength{\leftmarginii}{2.65em}
\setlength{\leftmarginiii}{3em}
\renewcommand{\labelenumi}{\textbf{\theenumi.}}
\pagestyle{empty}
\raggedbottom
\allowdisplaybreaks
\newenvironment{examproblems}[1]{%
  \begin{enumerate}%
  \setcounter{enumi}{#1}%
  \setlength{\topsep}{0.35em}%
  \setlength{\partopsep}{0pt}%
  \setlength{\itemsep}{0.48em}%
  \setlength{\parsep}{0.18em}%
}{\end{enumerate}}
\newenvironment{examsubparts}{%
  \begin{enumerate}%
  \setlength{\topsep}{0.18em}%
  \setlength{\partopsep}{0pt}%
  \setlength{\itemsep}{0.16em}%
  \setlength{\parsep}{0.1em}%
}{\end{enumerate}}
\newenvironment{examinstructions}{%
  \par\begingroup\itshape
}{%
  \par\endgroup\vspace{0.18em}
}
\makeatletter
\renewcommand\section{\@startsection{section}{1}{0pt}%
  {-1.25ex plus -.25ex minus -.1ex}%
  {0.55ex plus .1ex}%
  {\normalfont\large\bfseries}}
\renewcommand{\@maketitle}{%
  \newpage\null\vskip -1.2em%
  {\footnotesize\bfseries
    \href{https://www.ufl.edu/}{University of Florida}, \href{https://math.ufl.edu/}{Department of Mathematics}\par}%
  \vskip 0.18em%
  \tagstructbegin{tag=Title}%
    {\LARGE\bfseries\href{https://gma.math.ufl.edu/exams/probability-phd/}{Probability PhD Exam}, May 2023\par}%
  \tagstructend%
  \vskip 0.35em%
  \tagmcbegin{artifact}%
    \hrule height 1pt%
  \tagmcend%
  \vskip 0.55em%
}
\makeatother
\title{Probability PhD Exam, May 2023}
\author{}
\date{}
\begin{document}
\maketitle
\thispagestyle{empty}
\begin{examinstructions}
Carefully justify your answers.
\end{examinstructions}
\begin{examproblems}{0}
\item
Give the mathematical definition of:
\begin{examsubparts}
\item[\textnormal{1.}]
conditional expectation,
\item[\textnormal{2.}]
martingale,
\item[\textnormal{3.}]
Poisson process,
\item[\textnormal{4.}]
standard Brownian motion.
\end{examsubparts}
\item
This problem has three parts.
\begin{examsubparts}
\item[\textnormal{1.}]
State the weak and strong law of large numbers for i.i.d. random variables.
\item[\textnormal{2.}]
Prove the strong law of large numbers under the additional assumption of finite fourth moment \(\bigl(\mathbb{E}[X_1^4]<+\infty\bigr)\).
\item[\textnormal{3.}]
Let \(\{B_t\}\) be a standard Brownian motion. Prove that \(\left\{\frac{B_n}{n}\right\}\) converges to~\(0\) almost surely as \(n\to+\infty\) (\(n\in\mathbb{N}\)).
\end{examsubparts}
\item
This problem has two parts.
\begin{examsubparts}
\item[\textnormal{1.}]
Let~\(X\) be a random variable with a standard Gaussian distribution. Find the probability density function of \(e^X\).
\item[\textnormal{2.}]
Let~\(X\) be a random variable with a standard Cauchy distribution. Find the probability density function of \(\frac{1}{X}\).
\end{examsubparts}
\item
This problem has two parts.
\begin{examsubparts}
\item[\textnormal{1.}]
Let \(p\ge 1\). Prove that if~\(X\) is a random variable, then 
\[
\mathbb{E}[|X|^p]=\int_0^{+\infty}p x^{p-1}\mathbb{P}(|X|\ge x)\,dx.
\]
\item[\textnormal{2.}]
Let \(p\ge 1\). Let~\(X\) and~\(Y\) be two independent random variables with \(\mathbb{E}[Y]=0\). Show that 
\[
\mathbb{E}[|X+Y|^p]\ge \mathbb{E}[|X|^p].
\]
\end{examsubparts}
\item
Let \(\{a_n\}_{n\ge1}\) be a sequence of real numbers such that \(0<\sum_{k=1}^{+\infty}|a_k|^2<+\infty\). Denote 
\[
\|a_n\|=\sqrt{\sum_{k=1}^n|a_k|^2}
\quad\text{and}\quad
\|a\|=\sqrt{\sum_{k=1}^{+\infty}|a_k|^2}.
\]
 Let \(\{X_n\}_{n\ge1}\) be a sequence of i.i.d. symmetric Bernoulli random variables, that is, 
\[
\forall k\ge1,\qquad \mathbb{P}(X_k=1)=\mathbb{P}(X_k=-1)=\frac12.
\]
 For \(n\ge1\), denote \(S_n=\sum_{k=1}^n a_kX_k\), and denote \(\mathcal{F}_n=\sigma\{S_1,\ldots,S_n\}\).
\begin{examsubparts}
\item[\textnormal{1.}]
Prove that \(\{S_n\}\) is an \(\{\mathcal{F}_n\}\)-martingale.
\item[\textnormal{2.}]
Prove that \(\{S_n\}\) converges almost surely.
\item[\textnormal{3.}]
Prove that for all \(n\ge1\), for all \(\lambda\in\mathbb{R}\), 
\[
\mathbb{E}[e^{\lambda S_n}]\le e^{\frac{\lambda^2\|a_n\|^2}{2}}.
\]
 (Hint: \(\frac{e^x+e^{-x}}{2}\le e^{x^2/2}\))
\item[\textnormal{4.}]
Use question 3 and the symmetry of \(S_n\) to prove that for all \(n\ge1\), for all \(x\ge0\), 
\[
\mathbb{P}(|S_n|\ge x)\le 2e^{-\frac{x^2}{2\|a_n\|^2}}.
\]
 (Hint: The minimum of the function \(\lambda\mapsto e^{-\lambda x}e^{\frac{\lambda^2\|a_n\|^2}{2}}\) is attained at \(\lambda=\frac{x}{\|a_n\|^2}\).)
\item[\textnormal{5.}]
Deduce the Khintchine inequality: \(\forall p\ge1\), \(\forall n\ge1\), 
\[
\mathbb{E}\left[\left|\sum_{k=1}^n a_kX_k\right|^p\right]^{\frac1p}\le C\|a_n\|,
\]
 where~\(C\) is a numerical constant depending on~\(p\) only (you do not need to compute~\(C\)).

\par
(Hint: \(\mathbb{E}[|X|^p]=\int_0^{+\infty}p x^{p-1}\mathbb{P}(|X|\ge x)\,dx\).)
\end{examsubparts}
\item
This problem has three parts.
\begin{examsubparts}
\item[\textnormal{1.}]
Give the definition of convergence of a sequence of random variables in probability and in distribution.
\item[\textnormal{2.}]
Which mode of convergence is stronger between convergence in probability and in distribution? Prove it.
\item[\textnormal{3.}]
Give a counterexample showing that convergence in probability is not equivalent to convergence in distribution.
\end{examsubparts}
\item
Let \(\{B_t\}\) be a standard Brownian motion. Define, for \(x\in\mathbb{R}\), 
\[
T_x=\min\{t\ge0:B_t=x\}.
\]
\begin{examsubparts}
\item[\textnormal{1.}]
Prove that for all \(x\in\mathbb{R}\), \(T_x\) is finite almost surely (that is, \(\mathbb{P}(T_x<+\infty)=1\)).
\item[\textnormal{2.}]
Find \(\mathbb{P}(T_{-2}<T_1)\). You can use a picture as justification.
\end{examsubparts}
\item
Let \(n\in\mathbb{N}\). Let \(X_1,\ldots,X_n\) be i.i.d. random variables in \(L_1\). Find \(\mathbb{E}[X_1\mid X_1+\cdots+X_n]\).
\end{examproblems}
\end{document}
